\documentclass[10pt,a4paper]{amsart}
\usepackage[margin=19mm]{geometry}
\usepackage{amsmath,amssymb,mathtools}
\usepackage[scr=rsfs,cal=euler]{mathalfa}
\usepackage{xfrac}
\usepackage[unicode,hidelinks,bookmarks=false]{hyperref}
\newcommand{\ie}{i.e.\ }
\newcommand{\cf}{cf.\ }
\newcommand{\resp}{respectively\ }
\newcommand{\Subsection}{\subsection}
\providecommand{\nobreakdash}{\nobreak\hbox{-}}
\setlength{\emergencystretch}{2em}
\multlinegap0pt
\usepackage{amssymb}
\newtheorem{prop}{}
\title{\vspace{-2cm} {\footnotesize \textnormal{Published Open Access in \textit{Journal of Differential Equations}. \href{https://doi.org/10.1016/j.jde.2026.114334}{Download the official version here.}}} \\ [1ex] Nonlocal pseudosymmetries and Bäcklund transformations as $\mathcal{C}$-morphisms}
\author{Diego Catalano Ferraioli, Tarc\'isio Castro Silva}
\date{Source: arXiv:2506.01932v6 (CC BY 4.0)}
\begin{document}
\maketitle

\noindent\emph{Source and licence.} \vspace{-2cm} {\footnotesize \textnormal{Published Open Access in \textit{Journal of Differential Equations}. \href{https://doi.org/10.1016/j.jde.2026.114334}{Download the official version here.}}} \\ [1ex] Nonlocal pseudosymmetries and Bäcklund transformations as $\mathcal{C}$-morphisms. Version arXiv:2506.01932v6 (CC BY 4.0), distributed by arXiv under the Creative Commons Attribution 4.0 International licence, http://creativecommons.org/licenses/by/4.0/, as recorded in the arXiv licence metadata for this exact version and shown on the abstract page https://arxiv.org/abs/2506.01932v6. Full licence notice: \texttt{licenses/arxiv-2506.01932-CC-BY-4.0.txt}.\par\smallskip

\noindent\textit{Editorial scope.} Counted unit: the article's prop prop:pseudoE (source0.tex lines 1563--1624). The article's complete original proof of it is delivered with it (source0.tex lines 1626--1717, one \texttt{proof} environment of 92 lines). No mathematics is written, repaired or re-derived: the statement and the proof are the article's own lines, in the article's own notation, and no retained line is substituted or re-flowed.  No further source-local context is needed: every cross-reference of every retained line resolves inside the two delivered spans.  Nothing was omitted from any retained span, so the package carries no omission record.  No retained line contains a citation, so the wrapper carries no bibliography and no bibliography entry.  No retained line contains a figure environment, an image-inclusion command or drawing code, so no image is delivered and the package declares no asset dependency.  Editorial formatting only, in addition: an AMS \texttt{amsart} preamble that restates the article's own declarations and macros used by the retained lines, namely the environments \texttt{prop} on separate counters, together with the standard packages those lines need.  The retained environments print in this document as Proposition 1 in order of appearance, whereas the article prints the same items with the numbering of its own full document.  The complete native source of the article as distributed in the arXiv e-print for this exact version is bundled unchanged as \texttt{sources/179180/source0.tex}.
\begin{prop}
%%LEAP%%%\label{prop11}
\label{prop:pseudoE}%
Let $\mathcal{E}=\{\mathbf{F=0}\}\subset J^{k}(\pi )$, with
$\text{$\mathbf{F}$}=(F^{1},...,F^{p})$, be a formally integrable
differential equation and $\mu =U_{1}dx^{1}+...+U_{n}dx^{n}$ an horizontally
closed $1$-form on $\mathcal{E}^{(\infty )}$, i.e., such that
%
%e16 #&#
\begin{equation}
D_{j}\,U_{i}-D_{i}\,U_{j}=0\mod{\mathcal{E}^{(\infty )}}.
\label{eq:U_psym}
%%LEAP%%%\label{eq16}
\end{equation}
%
Then, the vector field on $J^{\infty}(\pi )$
%
%e17 #&#
\begin{equation}
Y=\sum _{i}a_{i}\partial _{x_{i}}+\sum _{|\sigma |\geq 0}\sum _{j}b_{
\sigma}^{j}\partial _{u_{\sigma}^{j}}=\sum _{i}a_{i}D_{i}+\sum _{|
\sigma |\geq 0}\sum _{j}\left (D+U\right )_{\sigma}(\varphi ^{j})
\partial _{u_{\sigma}^{j}},
\label{eq:Y_pseudo_int}
%%LEAP%%%\label{eq17}
\end{equation}
%
where
%
\begin{equation*}
%
\begin{array}{l}
\varphi ^{j}:=Y\,\lrcorner \,\omega _{0}^{j},\qquad \qquad \omega _{0}^{j}=du^{j}-{
\displaystyle \sum _{i}}u_{i}^{j}dx_{i},
\vspace{10pt}
\\
b_{\sigma}^{j}:={\displaystyle \sum _{i}a_{i}\,u_{\sigma +1_{i}}^{j}+
\left (D+U\right )_{\sigma}(\varphi ^{j})},
\vspace{10pt}
\\
\left (D+U\right )_{\sigma}:=\left (D_{1}+U_{1}\right )^{\sigma _{1}}
\circ ...\circ \left (D_{n}+U_{n}\right )^{\sigma _{n}},
\vspace{10pt}
\end{array}
%
\end{equation*}
%
restricts on $\mathcal{E}^{(\infty )}$ to a pseudosymmetry $\bar{Y}$ of
$\mathcal{E}$ whenever $Y(F^{h})=0\mod{\mathcal{E}^{(\infty )}}$, i.e.,
%
%e18 #&#
\begin{equation}
\sum _{|\sigma |\geq 0}\sum _{j}\partial _{u_{\sigma}^{j}}\left (F^{h}
\right )\left (D+U\right )_{\sigma}(\varphi ^{j})=0
\mod{\mathcal{E}^{(\infty )}},
\label{eq:Y_tang}
%%LEAP%%%\label{eq18}
\end{equation}
%
for any $h=1,...,p$, which is equivalent to the tangency condition of
$Y$ to $\mathcal{E}^{(\infty )}$.
\end{prop}

\begin{proof}
In view of t{(\ref{eq:Y_pseudo_int})} one gets (for ease of notation, we use
here the Einstein's summation convention)
%
\begin{equation*}
%
\begin{array}{l@{\,}l}
\left [D_{s},Y\right ] & =\left [D_{s},a_{i}D_{i}\right ]+\left [D_{s},
\left (D+U\right )_{\sigma}(\varphi ^{j})\partial _{u_{\sigma}^{j}}
\right ]
\vspace{5pt}
\\
& =D_{s}\left (a_{i}\right )\,D_{i}+\left [D_{s}+U_{s},\left (D+U
\right )_{\sigma}(\varphi ^{j})\partial _{u_{\sigma}^{j}}\right ]-
\left [U_{s},\left (D+U\right )_{\sigma}(\varphi ^{j})\partial _{u_{
\sigma}^{j}}\right ].
\end{array}
%
\end{equation*}
%
On the other hand, by t{(\ref{eq:U_psym})} one has
%
\begin{equation*}
\left (D_{s}+U_{s}\right )\circ \left (D+U\right )_{\sigma}=\left (D+U
\right )_{\sigma +1_{s}}\mod{\mathcal{E}^{(\infty )}}.
\end{equation*}
%
Hence, above identity can be rewritten as
%
\begin{equation*}
%
\begin{array}{l@{\,}l}
\left [D_{s},Y\right ] & =D_{s}(a_{i})D_{i}+\left (D+U\right )_{
\sigma +1_{s}}(\varphi ^{j})\partial _{u_{\sigma}^{j}}-\left (D+U
\right )_{\sigma}(\varphi ^{j})\partial _{u_{\sigma -1_{s}}^{j}}
\vspace{5pt}
\\
& \qquad -\left (D+U\right )_{\sigma}(\varphi ^{j})\partial _{u_{
\sigma}^{j}}\circ U_{s}+\left (D+U\right )_{\sigma}(\varphi ^{j})
\partial _{u_{\sigma}^{j}}\circ U_{s}
\vspace{5pt}
\\
& \qquad -U_{s}\circ (Y-a_{i}D_{i})\mod{\mathcal{E}^{(\infty )}},
\end{array}
%
\end{equation*}
%
which reduces to
%
%e19 #&#
\begin{equation}
[D_{s},Y]=\left (D_{s}a_{i}+a_{i}U_{s}\right )D_{i}-U_{s}Y
\mod{\mathcal{E}^{(\infty )}},
\label{eq:aux2}
%%LEAP%%%\label{eq19}
\end{equation}
%
in view of arbitrariness of $\sigma $. On the other hand t{(\ref{eq:aux2})}
entails that
%
\begin{equation*}
%
\begin{array}{l@{\,}l}
D_{s}\left (Y(D_{\sigma}F^{h})\right )-Y\left (D_{s}(D_{\sigma}F^{h})
\right ) & =\left (D_{s}a_{i}+a_{i}U_{s}\right )\,D_{i}\left (D_{
\sigma}F^{h}\right )-U_{s}\,Y(D_{\sigma}F^{h})
\vspace{5pt}
\\
& =-U_{s}\,Y(D_{\sigma}F^{h})\mod{\mathcal{E}^{(\infty )}},
\end{array}
%
\end{equation*}
%
i.e.,
%
%e20 #&#
\begin{equation}
Y\left (D_{\sigma +1_{s}}F^{h}\right )=U_{s}\,Y(D_{\sigma}F^{h})+D_{s}
\left (Y(D_{\sigma}F^{h})\right )\mod{\mathcal{E}^{(\infty )}},
\label{aux-1}
%%LEAP%%%\label{eq20}
\end{equation}
%
for any $\sigma $ and for any $h$. Thus, by repeatedly using t{(\ref{aux-1})}
with increasing $|\sigma |$, conditions
$Y(F^{h})=0\mod{\mathcal{E}^{(\infty )}}$, $h=1,...,p$, guarantee that
$Y\left (D_{\sigma}F^{h}\right )=0\mod{\mathcal{E}^{(\infty )}}$, i.e.,
that $Y$ is tangent to $\mathcal{E}^{(\infty )}$. Therefore, in view of
t{(\ref{eq:aux2})}, one gets that $Y$ restricts on
$\mathcal{E}^{(\infty )}$ to a pseudosymmetry $\bar{Y}$ of
$\mathcal{C}(\mathcal{E})$.
\end{proof}

\end{document}
